\section{Discussion}

\subsection{Average WER hides interactive tail risk}

Corpus-average \wer{} is necessary but insufficient for an interactive text-entry system. The
\texttt{large-v3} experiments illustrate why. A rare runaway continuation may affect only a
small number of utterances and therefore move the mean modestly, yet a single event can insert a
large block of unspoken text into the active document.

This suggests that interactive ASR evaluation should report error decomposition, repeated-run
behavior, and explicit catastrophic-event counts alongside mean \wer{}. The relevant engineering
question is not only how many reference words are wrong, but how the failure appears in the
editing surface a user must repair.

\subsection{Larger recognizers are not monotonically safer}

The measured checkpoint ladder resists a simple larger-is-better rule. Larger checkpoints reduce
some recognition errors, but the \texttt{large-v3} condition introduces a continuation failure
that is not observed in the same way on the smaller measured checkpoints. Likewise,
previous-text conditioning changes from beneficial, to neutral, to associated with a harmful
tail mode across the tested ladder.

Model choice in an interactive system therefore needs a broader objective than average accuracy.
Latency, memory, repetition behavior, and correction cost can matter as much as another fraction
of a WER point.

\subsection{Configuration is model-specific}

Beam width and previous-text conditioning both show that decoder guidance should be tied to a
specific model and workload. The campaign does not support globally disabling context or globally
maximizing beam width. Product defaults should remain evidence-linked to the checkpoint they are
intended to serve.

\subsection{Cross-platform has multiple layers}

Cross-platform claims should distinguish dependency resolution, numerical decode behavior,
activation/capture support, text injection, and recovery behavior. The present work strengthens
the decode and installability layers but does not collapse them into full end-to-end
certification. That stronger validation remains a separate experiment.

\subsection{Negative results improve the system}

The onset-padding, streaming, and centroid-merge experiments share a useful pattern. Each tested
an intuitive mechanism that could easily have been defended from first principles alone.
Measurement either narrowed the claim or rejected the mechanism.

This is especially valuable in a user-facing system because an apparently harmless feature can
consume CPU, complicate privacy guarantees, or create new failure modes. Reporting negative
results makes the relationship between the product and its evidence more auditable.

\subsection{Provenance is part of the experimental method}

The \texttt{large-v3} investigation began because a repeated benchmark disagreed with the
existing result and the repository preserved both artifacts. If a conventional latest-result
workflow had silently replaced the first run, the disagreement might have been interpreted as
ordinary benchmark drift rather than a reproducibility question.

For this project, provenance is therefore not only documentation. It changes which scientific
questions become visible.
